\documentclass{article}
\usepackage{amsmath}
\usepackage{amssymb}
\usepackage{graphicx}
\usepackage{booktabs}
\usepackage{array}
\usepackage{hyperref}
\usepackage{enumitem}
\usepackage{tikz}
\usepackage{float}
\usepackage{verbatim}

\title{DSC 257R: Unsupervised Learning \\ Homework 8}
\author{}
\date{}

\begin{document}

\maketitle

\section*{Experiments with the EM Algorithm and Kernel Density Estimation}

\begin{enumerate}
\item \textbf{Experiments with the EM algorithm.} In this problem, we'll generate data from a mixture of Gaussians and then see whether the EM algorithm is able to recover the correct clustering and the correct parameters of the mixture model.

\begin{enumerate}
\item Consider a mixture of four Gaussians in $\mathbb{R}^2$:
\[
p(x) = \pi_1 N(\mu_1, \Sigma_1) + \pi_2 N(\mu_2, \Sigma_2) + \pi_3 N(\mu_3, \Sigma_3) + \pi_4 N(\mu_4, \Sigma_4)
\]
where the mixing weights are
\[
\pi_1 = 0.1, \pi_2 = 0.2, \pi_3 = 0.5, \pi_4 = 0.2
\]
the cluster means are
\[
\mu_1 = \begin{bmatrix}7 \\ 6\end{bmatrix}, \mu_2 = \begin{bmatrix}0 \\ 0\end{bmatrix}, \mu_3 = \begin{bmatrix}-1 \\ 4\end{bmatrix}, \mu_4 = \begin{bmatrix}5 \\ -2\end{bmatrix}
\]
and the covariance matrices are
\[
\Sigma_1 = \begin{bmatrix}1 & 0 \\ 0 & 25\end{bmatrix}, \Sigma_2 = \begin{bmatrix}1 & 0 \\ 0 & 1\end{bmatrix}, \Sigma_3 = \begin{bmatrix}9 & 1 \\ 1 & 1\end{bmatrix}, \Sigma_4 = \begin{bmatrix}16 & -6 \\ -6 & 4\end{bmatrix}
\]

Write a routine that generates 400 points at random from this mixture model. Plot the points, using colors to distinguish the four components. Put this plot in your solution.

\item Write a routine that takes a two-dimensional Gaussian mixture (means and covariance matrices) and plots ellipses corresponding to each component. Using this routine, plot all the data points (from (a)) in the same color and then draw four ellipses on the same plot, depicting representative contours of the four Gaussians.

\item Look at the plots from (a) and (b) and check that the shapes of the clusters correspond to what you would expect from the covariance matrices. One of the clusters is small and is partially contained within another cluster. Which clusters are these?

\item Now run the EM algorithm on the data. Plot the data, and using the procedure from (b), draw four ellipses on the same plot, depicting the four Gaussians returned by EM. How close are they to the correct Gaussians?

\item Plot the data again, but this time color each point according to its most likely mixture component. How does this clustering compare to the true component memberships (from (a))? Are the discrepancies understandable?

\item Finally, let's see how EM's solution evolves over iterations of local search. Show four plots, depicting the progress of EM at different iterations.
\end{enumerate}

\item \textbf{Kernel density estimation.} Suppose we have samples $x_1, x_2, \ldots, x_n \in \mathbb{R}^d$ from an unknown distribution $f$, and we want to come up with a model of $f$. One simple way to do this is using the kernel density estimator $\widehat{f}$, which approximates $f$ using a mixture of $n$ Gaussians, each with weight $1/n$, and each centered at one of the data points:
\[
\widehat{f}: \frac{1}{n} N(x_1, \sigma^2 I_d) + \frac{1}{n} N(x_2, \sigma^2 I_d) + \cdots + \frac{1}{n} N(x_n, \sigma^2 I_d)
\]
Here $\sigma$ must be selected carefully. It is usually called the bandwidth parameter.

The file \texttt{samples.txt} contains 1,000 samples drawn from $f$ (using rejection sampling).
\begin{comment}
\begin{center}
\includegraphics[width=0.7\textwidth]{density_plot.png}
\captionof{figure}{Density $f(x)$}
\end{center}
\end{comment}

\begin{enumerate}
\item Plot a histogram of the 1,000 samples using 20 bins.

\item Now apply kernel density estimation using just the first 100 samples in \texttt{samples.txt}, that is, $n=100$. For the bandwidth parameter, try the following settings: $0.25, 0.5, 0.75, 1.0$.
\begin{itemize}
\item For each of the four bandwidth settings, show a plot of the resulting density $\widehat{f}$.
\item What changes as the bandwidth increases? Give a one- or two-sentence description.
\item In this case, which of the four bandwidth settings yields a model $\widehat{f}$ that is closest to the original $f$, in your opinion?
\end{itemize}

\item Now show the plots of the kernel density estimate using all 1000 samples in \texttt{samples.txt}, that is, $n=1000$. Use the same four settings of the bandwidth.
\end{enumerate}

\end{enumerate}

\end{document}
